\documentclass[11pt,a4paper]{article}
\usepackage[utf8]{inputenc}
\usepackage{amsmath, amssymb, amsthm, hyperref}

\title{Formal Verification Framework for Erd\H{o}s Problem 265}
\author{Lean 4 Formalization Project}
\date{\today}

\begin{document}

\maketitle

\begin{abstract}
We present a formal verification framework in Lean 4 for the analysis of Erd\H{o}s Problem 265. This problem investigates the growth rate of integer sequences $a_n$ for which both $\sum 1/a_n$ and $\sum 1/(a_n - 1)$ are rational. We provide formal proofs for the single-sum growth limit, the rational lattice gap, and the stall-length bound which limits the impact of wild oscillators. Furthermore, we establish a structural roadmap for establishing the $\beta \le 2$ growth ceiling and the $\beta \le 3$ algebraic envelope. All verified results are formalized with zero axioms and zero \textit{sorry} statements.
\end{abstract}

\section{Introduction}
Erd\H{o}s Problem 265 concerns the tension between the additive structure of reciprocals and the multiplicative growth of sequence terms. The problem asks for the supremum of the double-exponential base $\beta = \limsup a_n^{1/2^n}$ for sequences satisfying the dual-rationality constraint. This project establishes a formal roadmap for mapping the boundaries of this growth using the Lean 4 interactive theorem prover.

\section{Verified Structural Boundaries}
The following properties have been established as absolute Diophantine limits:
\begin{enumerate}
    \item \textbf{The Folklore Theorem:} For sequences with a uniform growth limit, the double-exponential base is strictly bounded by 2.
    \item \textbf{The Lattice Gap:} For any rational target $p/q$, the tail sum residual is restricted to a discrete lattice with minimum spacing $1/(q P_N)$, where $P_N$ is the prefix product.
    \item \textbf{The Stall Length Bound:} If a sequence attempts a massive growth jump (a ``wild oscillator''), the discrepancy between the step size and the rational lattice forces the sequence into a stall phase of strictly bounded length.
\end{enumerate}

\section{The Resolution Frontier}
While the framework provides a rigorous foundation, the absolute resolution of the problem depends on establishing the link between discrete residuals and continuous growth rates. We formalize these dependencies as structured hypotheses:
\begin{enumerate}
    \item \textbf{The $\beta \le 3$ Envelope:} Proved conditional on the ``Analytic Squeeze'' link between tail sums and prefix products, exploiting a 2-adic valuation ratchet.
    \item \textbf{The $\beta \le 2$ Barrier:} Established as a structural hypothesis, where the stall phase required by a wild oscillator triggers a prefix-inflation penalty that contradicts the trajectory.
    \item \textbf{Supremum Existence:} Formalizing the convergence of sequences to specific rational targets at the $\beta = 2$ threshold remains an open analytic problem.
\end{enumerate}

\section{Conclusion}
This repository provides a machine-verified technical map of Erd\H{o}s Problem 265, identifying the exact algebraic and Diophantine limits where sequence growth is constrained by the rigidity of the integers.

\end{document}